\begin{theorem} \label{thm:q=5}
The group $G{\left(m, 5 \right)}$ is not the automorphism group of any digraph for any $m \geqslant 2$.
\end{theorem}
\begin{proof}
Let $G = G{\left(m, 5 \right)}$. It suffices to show that $G$ is not the automorphism group of any union of its four non-trivial orbital digraphs. We know already from Lemma~\ref{lm:Gamma_A,B} that $G$ is not the automorphism group of $\Gamma_A$ or $\Gamma_B$. Since $2^2 = 4 = -1$ in the field $F$, we have from Lemma~\ref{lm:Gamma_1,i} that $G$ is not the automorphism group of $\Gamma_1$ or $\Gamma_2$ either.

Consider the following orbital digraph union:
\begin{align*}
\Delta_A \cup \Delta_1 &= \left\{e_1 \otimes w,\ e_2 \otimes w,\ \left(e_1 \pm e_2 \right) \otimes w \mid w \in W \setminus \left\{0 \right\}\right\} \\
&= \left\{\left(1, 0 \right) \otimes w,\ \left(0, 1 \right) \otimes w,\ \left(1, 1 \right) \otimes w,\ \left(1, -1 \right) \otimes w \mid w \in W \setminus \left\{0 \right\}\right\}
\end{align*}
with respect to the basis $\left\{e_1, e_2 \right\}$ of $V$. Consider the image of this union under the linear transformation $\theta = \begin{bsmallmatrix} 1 & 1 \\ 1 & -1 \end{bsmallmatrix} \circ I$, where $I$ denotes the $m \times m$ identity matrix. If we take, for example, a vector of the form $x = \left(1, 1 \right) \otimes w$ for some $w \in W \setminus \left\{0 \right\}$, then $x \in \Delta_A \cup \Delta_1$ and its image under $\theta$ is given by
\begin{displaymath}
x^\theta = \left(\left(1, 1 \right) \begin{bmatrix} 1 & 1 \\ 1 & -1 \end{bmatrix} \right) \otimes \left(wI \right) = \left(2, 0 \right) \otimes w = \left(1, 0 \right) \otimes \left(2w \right) \in \Delta_A \cup \Delta_1.
\end{displaymath}
It is easy to check using a similar calculation that if $x$ is a vector of the form $\left(1, 0 \right) \otimes w$ or $\left(0, 1 \right) \otimes w$ or $\left(1, -1 \right) \otimes w$, then the image of $x$ under $\theta$ lies in $\Delta_A \cup \Delta_1$.

We have now seen that $\theta$ preserves the set $\Delta_A \cup \Delta_1$, which means $\theta$ is an automorphism of $\Gamma_A \cup \Gamma_1$ by Lemma~\ref{lm:endomorphism-automorphism}. Since $\theta$ fixes the zero vector but does not lie in $G_0 = D_8 \circ \GL{\left(m, 5 \right)}$, it cannot be an element of $G$. Therefore $G$ is not the full automorphism group of $\Gamma_A \cup \Gamma_1$.

Now observe that
\begin{align*}
\Delta_A \cup \Delta_2 &= \left\{e_1 \otimes w,\ e_2 \otimes w,\ \left(e_1 \pm 2 e_2 \right) \otimes w \mid w \in W \setminus \left\{0 \right\}\right\} \\
&= {\left\{\left(1, 0 \right) \otimes w,\ \left(0, 1 \right) \otimes w,\ \left(1, 2 \right) \otimes w,\ \left(1, -2 \right) \otimes w \mid w \in W \setminus \left\{0 \right\}\right\}}.
\end{align*}
It can be checked that the linear transformation $\begin{bsmallmatrix} 1 & 2 \\ 2 & 1 \end{bsmallmatrix} \circ I$ preserves the set $\Delta_A \cup \Delta_2$, hence it is an automorphism of $\Gamma_A \cup \Gamma_2$ by Lemma~\ref{lm:endomorphism-automorphism}. Once again, this is an automorphism that does not lie in $G$, so $G$ is not the full automorphism group of $\Gamma_A \cup \Gamma_2$.

Finally, we have
\begin{align*}
\Delta_1 \cup \Delta_2 &= \left\{\left(e_1 \pm e_2 \right) \otimes w,\ \left(e_1 \pm 2 e_2 \right) \otimes w \mid w \in W \setminus \left\{0 \right\}\right\} \\
&= {\left\{\left(1, 1 \right) \otimes w,\ \left(1, -1 \right) \otimes w,\ \left(1, 2 \right) \otimes w,\ \left(1, -2 \right) \otimes w \mid w \in W \setminus \left\{0 \right\}\right\}}.
\end{align*}
It can be checked that the linear transformation $\begin{bsmallmatrix} 1 & 0 \\ 0 & 2 \end{bsmallmatrix} \circ I$ fixes the set $\Delta_1 \cup \Delta_2$, hence it is an automorphism of $\Gamma_1 \cup \Gamma_2$ by Lemma~\ref{lm:endomorphism-automorphism}. This automorphism does not lie in $G$, so $G$ is not the full automorphism group of $\Gamma_1 \cup \Gamma_2$.

The remaining orbital digraph unions are $\Delta_A \cup \Delta_B$, $\Delta_B \cup \Delta_1$ and $\Delta_B \cup \Delta_2$, as well as each union of three non-trivial orbital digraphs. However, these are the complements of orbital digraph unions that we have already checked, and every digraph has the same automorphism group as its complement. Hence we conclude that the group $G{\left(m, 5 \right)}$ is not the automorphism group of any digraph, as it is not the automorphism group of any union of its non-trivial orbital digraphs.
\end{proof}

Next, take $q=7$, let $m \geqslant 2$ be an integer, and let $G = G{\left(m, 7 \right)}$. As we are now working in the field $F = \GF {\left(7 \right)}$, Lemma~\ref{lm:suborbit-equivalence} gives us $\Delta_2 = \Delta_3 = \Delta_4 = \Delta_5$ and $\Delta_1 = \Delta_6$. Therefore the orbitals of $G$ are those corresponding to the suborbits $\delta$, $\Delta_A$, $\Delta_B$, $\Delta_1$ and $\Delta_2$, so $G$ once again has rank $5$.

\begin{theorem} \label{thm:q=7}
The group $G{\left(m, 7 \right)}$ is not the automorphism group of any digraph for any $m \geqslant 2$.
\end{theorem}
\begin{proof}
The orbital digraphs $\Gamma_A$, $\Gamma_B$ and $\Gamma_1$ do not have $G$ as their automorphism group according to Lemmas \ref{lm:Gamma_A,B} and \ref{lm:Gamma_1,i}.

The suborbit $\Delta_2$ as defined in Section \ref{sc:suborbits} is given by
\begin{align*}
\Delta_2 &= \left\{\left(e_1 \pm 2 e_2 \right) \otimes w,\ \left(e_1 \pm 2^{-1} e_2 \right) \otimes w \mid w \in W \setminus \left\{0 \right\}\right\} \\
&= \left\{\left(1, 2 \right) \otimes w,\ \left(1, 3 \right) \otimes w,\ \left(1, 4 \right) \otimes w,\ \left(1, 5 \right) \otimes w \mid w \in W \setminus \left\{0 \right\}\right\}.
\end{align*}
It can be checked that this set is preserved by the linear transformation $\begin{bsmallmatrix} 1 & 2 \\ 2 & 1 \end{bsmallmatrix} \circ I$. Therefore by Lemma~\ref{lm:endomorphism-automorphism} this transformation is an automorphism of $\Gamma_2$, and since it does not lie in $G$ we can conclude that $\Aut(\Gamma_2) \neq G$.

Similarly, observe that
\begin{displaymath}
\begin{array}{l@{}l@{}l}
\Delta_A \cup \Delta_1 &\ = \big\{&e_1 \otimes w,\ e_2 \otimes w,\ \left(e_1 \pm e_2 \right) \otimes w \mid w \in W \setminus {\left\{0 \right\}} \big\} \\
&\ = \big\{&\left(1, 0 \right) \otimes w,\ \left(0, 1 \right) \otimes w,\ \left(1, 1 \right) \otimes w,\ \left(1, 6 \right) \otimes w \mid w \in W \setminus {\left\{0 \right\}} \big\}, \\
\Delta_A \cup \Delta_2 &\ = \big\{&e_1 \otimes w,\ e_2 \otimes w,\ \left(e_1 \pm 2 e_2 \right) \otimes w,\ \left(e_1 \pm 2^{-1} e_2 \right) \otimes w \mid w \in W \setminus {\left\{0 \right\}} \big\} \\
&\ = \big\{&\left(1, 0 \right) \otimes w,\ \left(0, 1 \right) \otimes w,\ \left(1, 2 \right) \otimes w,\ \left(1, 3 \right) \otimes w,\ \left(1, 4 \right) \otimes w, \\
&&\left(1, 5 \right) \otimes w \mid w \in W \setminus {\left\{0 \right\}} \big\}, \\
\Delta_1 \cup \Delta_2 &\ = \big\{&\left(e_1 \pm e_2 \right) \otimes w,\ \left(e_1 \pm 2 e_2 \right) \otimes w,\ \left(e_1 \pm 2^{-1} e_2 \right) \otimes w \mid w \in W \setminus {\left\{0 \right\}} \big\} \\
&\ = \big\{&\left(1, \mu \right) \otimes w \mid \mu \in F \setminus {\left\{0 \right\}},\ w \in W \setminus {\left\{0 \right\}} \big\}.
\end{array}
\end{displaymath}
It can be checked that these sets are fixed by the linear transformations $\begin{bsmallmatrix} 1 & 1 \\ 1 & -1 \end{bsmallmatrix} \circ I$, $\begin{bsmallmatrix} 1 & 2 \\ 2 & 1 \end{bsmallmatrix} \circ I$ and $\begin{bsmallmatrix} 1 & 0 \\ 0 & 2 \end{bsmallmatrix} \circ I$, respectively. As these transformations do not lie in $G$, it follows from Lemma~\ref{lm:endomorphism-automorphism} that $G$ is not the automorphism group of $\Gamma_A \cup \Gamma_1$, nor $\Gamma_A \cup \Gamma_2$, nor $\Gamma_1 \cup \Gamma_2$.

Each of the remaining unions of non-trivial orbital digraphs is the complement of one of the unions we have already checked. Hence $G{\left(m, 7 \right)}$ is not the automorphism group of any digraph since it is not the full automorphism group of any union of its non-trivial orbital digraphs.
\end{proof}

We now consider the case where $q = 13$, so let $G = G{\left(m, 13 \right)}$ and observe from Lemma~\ref{lm:suborbit-equivalence} that
\begin{align*}
&\Delta_1 = \Delta_{12}, \\
&\Delta_2 = \Delta_6 = \Delta_7 = \Delta_{11}, \\
&\Delta_3 = \Delta_4 = \Delta_9 = \Delta_{10}, \\
&\Delta_5 = \Delta_8.
\end{align*}
Therefore the set
\begin{displaymath}
\mathcal{O} = \left\{\delta,\ \Delta_A,\ \Delta_B,\ \Delta_1,\ \Delta_2,\ \Delta_3,\ \Delta_5 \right\}
\end{displaymath}
is a complete set of distinct suborbits of $G$, so the rank of $G$ is $7$. The six non-trivial orbital digraphs are $\Gamma_A$, $\Gamma_B$, $\Gamma_1$, $\Gamma_2$, $\Gamma_3$ and $\Gamma_5$.

\begin{theorem} \label{thm:q=13}
The group $G{\left(m, 13 \right)}$ is not the automorphism group of any digraph for any $m \geqslant 2$.
\end{theorem}
\begin{proof}
First we examine the non-trivial orbital digraph unions that do not include $\Gamma_A$ or $\Gamma_B$. It can be checked that for each union of the orbital digraphs $\Gamma_1$, $\Gamma_2$, $\Gamma_3$ and $\Gamma_5$, the invertible linear transformation given in Table~\ref{tbl:q=13,endomorphisms} preserves the corresponding union of suborbits but does not lie in $G$. Each is an automorphism of the respective orbital digraph union by Lemma~\ref{lm:endomorphism-automorphism}.

\begin{table}
\centering
\begin{tabular}{cc||cc}
\hline
Orbital digraph union & Automorphism & Orbital digraph union & Automorphism \\
\hline
\vphantom{\Bigg(} $\Gamma_1$ & $\begin{bmatrix} 1 & 2 \\ 2 & 1 \end{bmatrix} \circ I$ & $\Gamma_2 \cup \Gamma_5$ & $\begin{bmatrix} 1 & 0 \\ 0 & 4 \end{bmatrix} \circ I$ \\
\vphantom{\Bigg(} $\Gamma_2$ & $\begin{bmatrix} 1 & 1 \\ 5 & -5 \end{bmatrix} \circ I$ & $\Gamma_3 \cup \Gamma_5$ & $\begin{bmatrix} 1 & 2 \\ 2 & 1 \end{bmatrix} \circ I$ \\
\vphantom{\Bigg(} $\Gamma_3$ & $\begin{bmatrix} 1 & 1 \\ 5 & -5 \end{bmatrix} \circ I$ & $\Gamma_1 \cup \Gamma_2 \cup \Gamma_3$ & $\begin{bmatrix} 1 & 0 \\ 0 & 2 \end{bmatrix} \circ I$ \\
\vphantom{\Bigg(} $\Gamma_5$ & $\begin{bmatrix} 1 & 1 \\ 1 & -1 \end{bmatrix} \circ I$ & $\Gamma_1 \cup \Gamma_2 \cup \Gamma_5$ & $\begin{bmatrix} 1 & 4 \\ 4 & -1 \end{bmatrix} \circ I$ \\
\vphantom{\Bigg(} $\Gamma_1 \cup \Gamma_2$ & $\begin{bmatrix} 1 & 4 \\ 4 & -1 \end{bmatrix} \circ I$ & $\Gamma_1 \cup \Gamma_3 \cup \Gamma_5$ & $\begin{bmatrix} 1 & 2 \\ 2 & 1 \end{bmatrix} \circ I$ \\
\vphantom{\Bigg(} $\Gamma_1 \cup \Gamma_3$ & $\begin{bmatrix} 1 & 0 \\ 0 & 4 \end{bmatrix} \circ I$ & $\Gamma_2 \cup \Gamma_3 \cup \Gamma_5$ & $\begin{bmatrix} 1 & 1 \\ 1 & -1 \end{bmatrix} \circ I$ \\
\vphantom{\Bigg(} $\Gamma_1 \cup \Gamma_5$ & $\begin{bmatrix} 1 & 0 \\ 0 & 5 \end{bmatrix} \circ I$ & $\Gamma_1 \cup \Gamma_2 \cup \Gamma_3 \cup \Gamma_5$ & $\begin{bmatrix} 1 & 0 \\ 0 & 2 \end{bmatrix} \circ I$ \\
\vphantom{\Bigg(} $\Gamma_2 \cup \Gamma_3$ & $\begin{bmatrix} 1 & 1 \\ 5 & -5 \end{bmatrix} \circ I$ \\
\hline
\end{tabular}
\caption{A linear automorphism not in $G{\left(m, 13 \right)}$ for each of the given unions of orbital digraphs. The $m \times m$ identity matrix is denoted by $I$.}
\label{tbl:q=13,endomorphisms}
\end{table}

Suppose now that $\Gamma = \Gamma_B \cup \Gamma'$, where $\Gamma'$ is some union of $\Gamma_1$, $\Gamma_2$, $\Gamma_3$ and $\Gamma_5$. We have just seen that there exists some $\theta \in \GL{\left(2, q \right)} \circ \GL{\left(m, q \right)}$ that is an automorphism of $\Gamma'$ but that does not lie in $G$. But since $\theta$ preserves the set $\Delta_B$ of all non-simple tensors, it must also be an automorphism of $\Gamma_B$ by Lemma~\ref{lm:endomorphism-automorphism}. It follows that $\theta \in \Aut{\left(\Gamma \right)}$.

Finally, if a union of non-trivial orbital digraphs of $G$ includes $\Gamma_A$, then its complement does not include $\Gamma_A$ and has therefore already been checked.

Since $G$ is not the full automorphism group of any union of its non-trivial orbital digraphs, we conclude that $G$ is not the automorphism group of any digraph.
\end{proof}
